

%%%%%%%%%%%%%%%%%%%%%%%%%
\section{Soergel Bimodules}
\label{sec-sbim}
%%%%%%%%%%%%%%%%%%%%%%%%%

%%%%%%%%%%%%%%%%%%%%%%%%%%%%%%%%%%%%%%%%%%%%
\subsection{Realizations of Coxeter systems} \label{subsec-realization}
%%%%%%%%%%%%%%%%%%%%%%%%%%%%%%%%%%%%%%%%%%%%

For both Soergel's construction of Soergel bimodules, and for our construction of a diagrammatic category by generators and relations, the starting point will be the data of a realization
of a Coxeter system.

\begin{defn} \label{defn:realization} Let $\Bbbk$ be a commutative ring. A \emph{realization} of $(W,S)$ over $\Bbbk$ is a free, finite rank $\Bbbk$-module $\hg$, together with subsets $\{ \a_s^\vee \; | \; s \in S\} \subset \hg$ and $\{ \a_s \; | \; s \in S\} \subset \hg^* = \Hom_{\Bbbk}(\hg,\Bbbk)$, satisfying:
\begin{enumerate}
\item $\langle \alpha_s^\vee, \alpha_s \rangle = 2$ for all $s \in S$;
\item the assignment $s(v) \define v - \langle v, \alpha_s\rangle \alpha_s^{\vee}$ for all $v \in \hg$ yields a representation of $W$;
\item the technical condition in \eqref{eq:technical-condition} is satisfied. (Its description requires some background.)
\end{enumerate}
We will often refer to $\hg$ as a realization, however the choice of $\{ \alpha_s^\vee \}$ and $\{\alpha_s \}$ is always implicit. \end{defn}


Given a realization over $\Bbbk$ and a homomorphism $\Bbbk \to \Bbbk'$ we obtain a realization $\Bbbk' \otimes_{\Bbbk} \hg$ over $\Bbbk'$ by base change. We call a realization
\emph{faithful} if the action of $W$ on $\hg$ (and hence the contragredient action on $\hg^*$) is faithful. Base change does not preserve faithfulness in general. For us, the ability to
perform base change is the more important property, so we must allow realizations which are not faithful. For instance, any realization of the dihedral group with $m_{st}=m<\infty$ is also
a realization of the dihedral group with $m_{st}=2m, 3m, \ldots$, and is also a realization of the infinite dihedral group.

We call a realization \emph{symmetric} if $\langle \alpha_s^\vee, \alpha_t \rangle = \langle \alpha_t^\vee, \alpha_s \rangle$ for all $s, t \in S$.

\begin{ex} \label{ex:realizations} Some examples of realizations that we have in mind are the following:
  \begin{enumerate}
  \item Let $(W,S)$ be any Coxeter system of finite rank. Let $\Bbbk = \RM$ and $\hg = \bigoplus_{s \in S} \RM \alpha_s^{\vee}$. Define elements $\{ \alpha_s \} \subset \hg^*$ by
\begin{equation} \label{eq:justcos}
\langle \alpha_t^\vee, \alpha_s \rangle = -2\cos(\pi/m_{st})
\end{equation}
(by convention $m_{ss} =1$ and $\pi/\infty = 0$). Then $\hg$ is a symmetric realization of $(W,S)$, called the \emph{geometric representation} (see \cite[\S 5.3]{Humphreys}). Note that the subset $\{ \alpha_s \} \subset  \hg^*$ is linearly independent if and only if $W$ is finite.

\item
More generally, given a real vector space $\hg$ with subsets $\{ \alpha_s^{\vee} \} \subset \hg$ and $\{ \alpha_s \} \subset \hg^*$ satisfying \eqref{eq:justcos} then $\hg$ is a realization of $(W,S)$. In  \cite[\S 2]{Soe3} Soergel builds his theory of Soergel bimodules for arbitrary Coxeter systems on a realization for which both $\{ \alpha_s^{\vee} \} \subset \hg$ and $\{ \alpha_s \} \subset \hg^*$ are linearly independent, and such that $\hg$ has minimal dimension with this property. To construct such a representation, Soergel mimics the construction of the action of an affine Weyl group on the Cartan subalgebra of an affine Kac-Moody group.

\item Let $(X, R, X^\vee, R^\vee)$ be a (reduced) root datum (see \cite[\S 7.4]{SprLAG} for notation) and let $\Delta \subset R$ be a set of simple roots. Let $(W,S)$ be the corresponding Weyl group and simple reflections. Then the triple $\hg = X$, $\{\alpha \; | \; \alpha \in \Delta \} \subset \hg$ and $\{\alpha^\vee \; | \; \alpha \in \Delta \} \subset \hg^* = X^\vee$ gives a faithful realization of $(W,S)$ over $\ZM$. We obtain a (potentially non-faithful) realization of $(W,S)$ over any $\Bbbk$ by extension of scalars.

\item
More generally, if $A$ is a generalized Cartan matrix and $\tg$ denotes the Cartan subalgebra of the corresponding Kac-Moody Lie algebra $\gg(A)$ (see \cite[Chapters 1 and 3]{KacBook}) then any choice of $\ZM$-lattices $\hg \subset \tg$ such that $\hg$ contains the root lattice and its dual lattice $\hg^* \subset \tg^*$ contains the coroot lattice yields a realization of the Weyl group $(W,S)$ of $\gg(A)$. In this way one obtains realizations over $\ZM$ (and hence over any $\Bbbk$) of any Coxeter system for which $m_{st} \in \{ 2,3,4,6, \infty\}$ for all $s \ne t \in S$. (Such Coxeter systems are called \emph{crystallographic}.)

\item Let $(W,S)$ be a Coxeter system for  which $m_{st} \in \{ 2,3,5,\infty\}$ and let $\Bbbk = \ZM[\phi]$ where $\phi = (1 + \sqrt{5})/2$ denotes the golden ratio. Let $\hg = \bigoplus_{s \in S} \Bbbk \alpha_s^{\vee}$ and define $\alpha_s \in \hg^*$ via
\[
\langle \alpha_s^\vee, \alpha_t \rangle = \begin{cases} 2 & \text{if $s = t$,} \\ 0 & \text{if $m_{st} =2 ,$} \\ -1 & \text{if $m_{st} =3 ,$} \\ -\phi & \text{if $m_{st} =5,$} \\ -2 & \text{if $m_{st} =\infty.$} \end{cases}
\]
Using that $-2\cos(\pi/5) = -\phi$ it follows from the example of the geometric realization that $(\hg, \{ \alpha_s^{\vee} \}, \{ \alpha_s \})$ is a (symmetric) realization of $(W,S)$ over $\Bbbk$. In particular, the finite reflection groups of types $H_3$ and $H_4$ have symmetric realizations over (any extension of) $\Bbbk$.

\item Let $(W,S)$ be the affine Weyl group of type $\widetilde{A_n}$ for $n \ge 2$. Let $\Bbbk = \ZM[q,q^{-1}]$ and $\hg = \bigoplus_{s \in S} \ZM[q,q^{-1}] \alpha_s^{\vee}$, and let the values of
$\langle \alpha_s^\vee, \alpha_t \rangle$ be encoded (as will be described soon) in the following matrix: \[A = \left( \begin{array}{ccccc} 2 & -1 & 0 & 0 & -q^{-1} \\ -1 & 2 & -1 & 0 & 0
\\ 0 & -1 & 2 & -1 & 0 \\ 0 & 0 & -1 & 2 & -q \\ -q & 0 & 0 & -q^{-1} & 2 \end{array} \right).\] 
(More precisely, this is the example when $n = 4$.)
This gives a realization of $W$. Specializing $q$ to an element of $\CM \setminus \RM$, one
obtains a realization of $W$ over $\CM$ which can not be obtained by extension of scalars from a realization defined over $\RM$.
\end{enumerate}
\end{ex}

Given a realization $(\hg, \{\alpha_s^{\vee} \}, \{\alpha_s\})$ of $(W,S)$ over $\Bbbk$ we can consider its \emph{Cartan matrix} $(\langle \alpha_s^\vee, \alpha_t \rangle ) _{s,t \in S}$.
Clearly a realization is symmetric if and only if its Cartan matrix is. Conversely, given a matrix $(a_{st})_{s,t \in S}$ such that $a_{ss}=2$, one can construct the free $\Bbbk$-module
$\hg = \bigoplus_{s \in S} \Bbbk \alpha_s^{\vee}$, and define $\alpha_s \in \hg^*$ by $\langle \alpha_s^\vee, \alpha_t \rangle = a_{st}$. When this yields a realization of $(W,S)$ we call
the matrix $(a_{st})$ a \emph{Cartan matrix} for $(W,S)$ over $\Bbbk$. Any realization for which $\{\a_s^\vee\}$ is a basis for $\hg$ can be reconstructed from its Cartan matrix; we call
such realizations \emph{minimal}.

\begin{ex} In Example \ref{ex:realizations} the realizations discussed in (1), (5) and (6) are minimal. The example in (3) is minimal if and only if the root system is simply connected and
of adjoint type (so that root lattice coincides with $\hg$). Examples (2) and (4) are not minimal in general. \end{ex}

\begin{remark} We expect that there is a rich Koszul duality theory for categories obtained from Soergel bimodules for arbitrary Coxeter systems (generalizing Soergel's description
\cite{Soe1} of the algebra of category $\OC$ in the case of Weyl groups). Here one expects Koszul duality to exchange $\hg$ and $\hg^*$, roots and coroots. In this setting it seems natural
to require both $\{ \alpha_s^{\vee} \} \subset \hg$ and $\{ \alpha_s \} \subset \hg^*$ to be linearly independent. This explains in part why we do not assume that our realizations are
minimal. \end{remark}

It is natural to ask under which conditions a matrix $(a_{st})_{s,t \in S}$ with $a_{ss} = 2$ is a Cartan matrix of $(W,S)$. A thorough discussion of this can be found in the appendix to
\cite{EDihedral}. We provide a short discussion here.

\begin{defn} \label{def:2qnum} Define the \emph{2-colored quantum numbers} $[k]_x$ and $[k]_y$ inside the ring $\ZM[x,y]$ inductively. One has $[0]_x = [0]_y = 0$, $[1]_x = [1]_y = 1$, and $[2]_x = x$, $[2]_y=y$. The other 2-colored quantum numbers are defined by the rules

\begin{subequations}
\begin{equation} [2]_x [k]_y = [k+1]_x + [k-1]_x, \end{equation}
\begin{equation} [2]_y [k]_x = [k+1]_y + [k-1]_y. \end{equation}
\end{subequations}

When $k$ is odd, $[k]_x = [k]_y$ and we shorten the notation to $[k]$. \end{defn}

Fix a pair $s,t \in S$ and let $x=a_{st}$ and $y=a_{ts}$. The condition that $(st)$ has order exactly $k$ is an algebraic condition on $x$ and $y$. For instance, when $\a_s$ and $\a_t$ are
linearly independent, the action of $(st)$ on their span has order $k$ when $[2k+1]=1$ and $[2k]_x = [2k]_y = 0$. In fact, this implies further that $2[k]_x = 2[k]_y = [2]_x [k]_y = [2]_y
[k]_x = 0$. This suggests that one should set \begin{equation} \label{eq:technical-condition} [m_{st}]_x = [m_{st}]_y = 0. \end{equation} This is the technical condition mentioned in
Definition \ref{defn:realization}. While \eqref{eq:technical-condition} is sufficient to imply that $(st)$ has order dividing $k$ on the span of the roots, it is independent of the
condition that $W$ acts on $\hg$. The reason that \eqref{eq:technical-condition} is required is to ensure that 2-colored Jones-Wenzl projectors are rotation-invariant, as discussed in
section \ref{subsec-JW}.

If either $x$ or $y$ is a non-zero-divisor, then \eqref{eq:technical-condition} is equivalent to the statement that $xy$ satisfies the minimal polynomial of the algebraic integer $4
\cos^2(\frac{\pi}{m})$. If this algebraic integer does not exist in $\Bbbk$ then $(W,S)$ does not admit a realization over $\Bbbk$. Any Coxeter system (of finite rank) admits a realization
over some ring of integers. Finally, we introduce one other technical condition.

\begin{defn} We call a realization \emph{balanced} if for every $s,t \in S$ one has $[m_{st}-1]_x = [m_{st}-1]_y = 1$. We refine this notion by calling the realization \emph{even-balanced}
(resp. \emph{odd-balanced}) if this property holds when $m_{st}$ is even (resp. odd). The opposite of even-balanced is \emph{even-unbalanced}. \end{defn}

The familiar Cartan matrices of Weyl groups are balanced. However, the Cartan matrix of type $A_2$ is not balanced when viewed as a realization of $G_2$. The exotic Cartan matrices for
type $\tilde{A}_n$, $n \ge 3$ given in Example (6) above are not balanced, except when $q=1$. Being balanced is equivalent to the existence of an unambiguous notion of positive roots in
$\hg^*$; when the realization is symmetric, being balanced is similar to the condition that simple roots form an obtuse angle. Faithful realizations are almost always even-balanced; any
minimal even-unbalanced realization over a domain is not faithful. Once again, a thorough discussion of these technicalities can be found in \cite{EDihedral}.

Fix a realization $(\hg, \{ \alpha_s^\vee \}, \{ \alpha_s \})$ of $(W,S)$ and let \[ R = \bigoplus_{m \ge 0} S^m(\hg^*) \] denote the symmetric algebra on $\hg^*$, which we view as a
graded $\Bbbk$-algebra with $\deg \hg^* = 2$. Then $W$ acts on $\hg^*$ via the contragredient representation ($s(\gamma) = \gamma - \langle \alpha_s^\vee, \gamma \rangle \alpha$ for all
$\g \in \hg^*$) and this extends to an action of $W$ on the algebra $R$ by graded automorphisms. We think of $R$ as the polynomial functions on $\hg$.

We let $R-\Mod$ and $R-\Bim$ denote the category of graded $R$-modules and bimodules respectively. We view $R-\Mod$ and $R-\Bim$ as graded categories; that is, as categories enriched in
graded $\Bbbk$-modules. We denote the grading shift by $(1)$: if $M = \bigoplus M^i$ is a graded (bi)module then $M(1)^i = M^{i+1}$. Degree $0$ maps of (graded) $R$-(bi)modules will be
denoted by $\Hom_0(M,N)$.


%%%%%%%%%%%%%%%%%%%%%%%%%%%%%%%%%%%%%%%%%%%%
\subsection{Technicalities}
%%%%%%%%%%%%%%%%%%%%%%%%%%%%%%%%%%%%%%%%%%%%

It is important to remember the key dichotomy in this paper: we will discuss two separate categories ``at once." Fix a realization of $(W,S)$, and consider the ring $R$ defined above. In
this chapter we will introduce Soergel's monoidal category $\SBim$, which is a full subcategory of $R-\Bim$. In chapter \ref{sec-diagsbim} we will define a diagrammatic category $\DC$ by
generators and relations, whose morphism spaces will be enriched in $R-\Bim$. One will need to make some assumptions on the realization in order for either category to ``behave well" (i.e.
in order for the Soergel Categorification Theorem to hold, and in order for double leaves to form a basis for Hom spaces; see the introduction). Whenever $\SBim$ behaves well, one can
construct an equivalence from $\DC$ to $\SBim$. However there are certain situations (for example when the characteristic of $\Bbbk$ is small, or when working over a complete local ring)
where the diagrammatic theory continues behaving well, but the bimodule theory either breaks down or has not yet been developed. In these cases, the diagrammatic theory seems to provide a
natural replacement for Soergel bimodules.\footnote{Another natural replacement is Fiebig's theory of sheaves on moment graphs.} This is one of the advantages of
the diagrammatic approach.

In this section we will discuss the technical assumptions one must make on the realization in order for the diagrammatics to behave well, and the further assumptions needed for Soergel
bimodules to behave well. The novice reader should ignore this section, and should stick with the geometric realization defined in part (1) of Example \ref{ex:realizations}. This section
may be overly pedantic, however, in view of current and future applications we make an effort to state all results in a natural level of generality.

Note that the very existence of a realization is already an assumption on the base ring $\Bbbk$: namely, that it contains certain algebraic integers.

\begin{assumption}(Demazure Surjectivity) \label{ass:Demazure Surjectivity} The map $\a_s \co \hg \to \Bbbk$ is surjective, for all $s \in S$. Evaluation at $\a_s^\vee$ gives a surjective
map $\hg^* \to \Bbbk$, for all $s \in S$. \end{assumption}

Whenever Demazure Surjectivity holds, there is some $\d \in \hg^*$ for which $\langle \a_s^\vee,\d \rangle = 1$. Moreover, $\a_s \ne 0$, so that $s(\d) = \d - \a_s \ne \d$.

If $2$ is invertible in $\Bbbk$ then Demazure Surjectivity holds, because $a_{ss}=2$. If $m_{st}$ is odd then both $\a_s$ and $\a_t$ (and $\a_s^\vee$ and $\a_t^\vee$) are surjective,
because the algebraic integer $4 \cos^2(\frac{\pi}{m_{st}})$ is invertible in any ring which contains it (see \cite{EDihedral}). Even when the ideal in $\Bbbk$ generated by $\langle
\a_t^\vee, \a_s \rangle$ as $t$ varies (for fixed $s$) is not the unit ideal, it is still possible that $\a_s$ is surjective when the realization is not minimal. Finally, the Demazure Surjectivity property is preserved by base change.

We will assume Demazure Surjectivity henceforth (with the exception of some remarks). In addition to standard ring-theoretic assumptions, this will be the only special assumption we need
to make in order for $\DC$ to be well-behaved.

Our arguments in Section \ref{subsec-SThmDiag} classifying the indecomposable objects in $\Kar(\DC)$ require that $\Bbbk$ is a complete local ring. This assumption is needed for either
category to satisfy the Krull-Schmidt theorem, and for idempotent lifting arguments to work. If $\Bbbk$ is not a complete local ring (for example $\ZM$) we have no idea how many
indecomposable Soergel bimodules there are, nor whether the Krull-Schmidt theorem holds. This is a typical situation in representation theory: one has a generic category (for example
representations of a finite group) defined over (some finite extension of) $\ZM$, but it is only after completing at a prime that one obtains a category in which it makes sense to discuss
indecomposable objects, do homological algebra etc. Moreover, the behavior at different primes can be vastly different.

We assume that $\Bbbk$ is a domain. In particular, it has no non-trivial idempotents, so that any graded $R$-(bi)module with $\End_0(M) = \Bbbk$ is indecomposable.

\begin{defn} Following Soergel \cite{Soe3}, we say that a realization $\hg$ over a field $\Bbbk$ is \emph{reflection faithful} if $\hg$ is a faithful representation of $W$, and if there is
a bijection between the set of reflections (i.e. the conjugates in $W$ of $S$) and the codimension one subspaces of $\hg$ that are fixed by some element of $W$. \end{defn}

This is a fairly serious assumption on a realization. For instance, no infinite Coxeter group admits a faithful representation over $\overline{\FM}_p$. Soergel constructs a reflection
faithful representation of any Coxeter group over $\RM$, using the approach mentioned in part (2) of Example \ref{ex:realizations}.

Soergel's theory gives techniques to study $\SBim$ defined for a reflection faithful representation over an infinite field $\Bbbk$ of characteristic $\ne 2$. Libedinsky \cite{Lib} has
shown that his results extend to the geometric realization as well. It seems plausible that many of Soergel's techniques could be adapted to other complete local rings and faithful
realizations over them, but the true generality of his results is unknown. We say that a realization is a \emph{Soergel realization} if it is faithful and Soergel's techniques can be
applied (i.e. if we can quote the Soergel Categorification Theorem).

Finally, while it will not affect the truth or falsehood of the Soergel Categorification Theorem, the assumption that the realization is balanced will drastically simplify both the study
of bimodules and the study of diagrammatics. We do not take this assumption in general.

%===========
\subsection{Demazure operators}
\label{subsec-polys}
%===========

Fix $s \in S$. We will extend the map $\langle \a_s^\vee, \cdot \rangle \co \hg^* \to \Bbbk$ to the \emph{Demazure operator} $\pa_s \co R \to R^s(-2)$, by the formula 
\[
\pa_s(f) = \frac{f - s(f)}{\a_s}.
\]
The numerator and denominator are both $s$-antiinvariant, so that the fraction, assuming it is well-defined, should lie in the subring $R^s$ of $s$-invariants.
Clearly this map makes sense for $f \in \hg^*$, and agrees with $\langle \a_s^\vee, \cdot \rangle$. Let us demonstrate that it makes sense in general.

Suppose that $\d$ is an element of $\hg^*$ such that $\langle \a_s^\vee, \d \rangle = 1$, guaranteed to exist by the assumption of Demazure Surjectivity. 

\begin{claim} Any element $f \in R$ can be written uniquely as $f = g\d + h$ for $g,h \in R^s$. \label{claim:unique decomposition}\end{claim}
	
\begin{proof} If $f$ is of this form then $f - s(f) = g(\d - s(\d)) = g \a_s$. The uniqueness of such an expression is now clear: if $g
\d + h = g' \d + h' = f$ then $f-s(f) = g \a_s = g' \a_s$. Since $\Bbbk$ is a domain $R$ is also, and therefore $g=g'$, so that $h=h'$. 

Clearly $a_{st} \d - \a_t \in \hg^*$ is $s$-invariant, lying in the kernel of $\langle \a_s^\vee, \cdot \rangle$. In particular, this implies that any polynomial in $R$ can
be expressed as a polynomial in $\d$ with coefficients in $R^s$. Moreover, $\d^2 = \d(\d + s(\d)) - \d s(\d)$, where both $\d + s(\d)$ and $\d s(\d)$ are $s$-invariant. Therefore any
polynomial in $\d$ can be written as $\d g + h$ for $g,h \in R^s$. \end{proof}

Thus defining $\pa_s(f) \define g$ makes sense, and agrees with the above formula. A similar argument shows that this alternative definition of $\pa_s$ is independent of
the choice of $\d$. The Demazure operator $\pa_s$ is a map $R^s$-bimodules, whose kernel is precisely $R^s$. It is easy to show that $\pa_s$ satisfies the \emph{twisted Leibniz rule}
\[\pa_s(fg) = a\pa_s(g) + \pa_s(a)(sg) \qquad \text{for all $f, g \in R$.} \] Unsurprisingly, Demazure Surjectivity implies that $\pa_s$ is surjective, for all $s \in S$.

Claim \ref{claim:unique decomposition} implies that $R$ is free of rank 2 over $R^s$, generated by $1$ and $\d$. In fact, the sets $\{1,\d\}$ and $\{-s(\d),1\}$ give dual bases of $R$ as
an $R^s$-module under the pairing $(g,h)=\pa_s(gh)$. This gives $R$ the structure of a graded Frobenius algebra over $R^s$. The element $\Delta_s = \d \ot 1 - 1 \ot s(\d) \in R \ot_{R^s}
R$ is independent of the choice of $\d$ with $\pa_s(\d)=1$, and satisfies $g \Delta_s = \Delta_s g$ for any $g \in R$.

\begin{remark} There is a unique choice $\d = \frac{\a_s}{2}$ such that the basis $\{1,\d\}$ is self-dual. It only exists when $2$ is invertible in $\Bbbk$. This was the choice used in the introduction. \end{remark}

\begin{remark} We have taken the assumption of Demazure Surjectivity precisely in order that $R$ would be a Frobenius extension of $R^s$, with trace map $\pa_s$. Without this assumption
the situation is less tractable. When the image of $\pa_s \co \hg^* \to \Bbbk$ is a non-trivial principal ideal, it will be true that $R$ is a Frobenius extension of $R^s$, but with
a rescaled trace map. If this image is a non-principal ideal, then $R$ is not even free over $R^s$. Other scenarios which our assumption forbids are $\pa_s = 0$ or $\a_s = 0$, as either
would imply $R=R^s$ (this is only a possibility in characteristic 2). \end{remark}

\begin{remark} If the realization is odd-balanced, then Demazure operators associated to $s \in S$ satisfy the braid relations. Otherwise, they do not. See \cite{EDihedral} for more
details. \end{remark}

%===========
\subsection{Soergel bimodules and standard bimodules}
\label{subsec-soergelandstandard}
%===========
In this section we give an introduction to Soergel bimodules and standard bimodules, following Soergel's ``classical'' treatment of the subject \cite{Soe3}.

For $s \in S$, let $B_s$ denote the $R$-bimodule $R \ot_{R^s} R (1)$, given by restriction followed by induction and a grading shift. Henceforth, $\ot$ will denote the tensor product over $R$, while $\ot_s$ will denote the tensor product over $R^s$.

Given a sequence $\un{w} = s_1s_2 \dots s_d$ the corresponding \emph{Bott-Samelson bimodule} is the tensor product
\[
B_{\un{w}} \define B_{s_1} \ot B_{s_2} \ot \ldots \ot B_{s_d}
\]
viewed as an $R$-bimodule under left and right multiplication. The Bott-Samelson bimodule $B_{\un{w}}$ is isomorphic to $R \ot_{s_1} R \ot_{s_2} R \ot \cdots \ot_{s_d} R (d)$. We let $\BSBim$ denote the full monoidal subcategory of $R-\Bim$ whose objects are Bott-Samelson bimodules (where as before, morphism spaces are graded $\Bbbk$-modules). Finally, we let $\SBim$ denote the idempotent closure or Karoubi envelope of (the additive, graded closure of) $\BSBim$, known as the category of \emph{Soergel bimodules}. That is, the indecomposable Soergel bimodules are the indecomposable direct summands of shifts of Bott-Samelson bimodules. Note that $\SBim$ is additive but not abelian.

There are also a number of other bimodules which play an important role in the theory. They are \emph{not} Soergel bimodules in general, because they do not appear as summands in Bott-Samelson
bimodules, only as submodules and quotients. These are the \emph{standard bimodules}. For $w \in W$, let $R_w$ denote the $R$-bimodule which is isomorphic to $R$ as a $\Bbbk$-module, and
where the left action of $f \in R$ is multiplication by $f$, while the right multiplication is multiplication by $w(f)$. It is clear that $R_w \ot R_v \cong R_{wv}$. We refer to the
additive monoidal category consisting of all direct sums of grading shifts of $R_w$ as $\StdBim$. This monoidal category is generated by $R_s$ for $s \in S$. A prototypical
object is $R_{\un{w}} \define R_{s_1} \ot R_{s_2} \ot \cdots R_{s_d}$. Unlike for Bott-Samelson bimodules, one has $R_{\un{w}} \cong  R_{\un{w}'}$ if $w = w'$.

It is useful to picture tensor products of bimodules $B_s$ and $R_w$, for example $B_s \ot B_t \ot R_w \ot B_s$, as being separators or dividers between regions, with regions corresponding
to the $\ot$ signs as well as to the left and right sides. A standard tensor in such a bimodule consists of a polynomial in each region. The bimodule encodes certain rules about how
polynomials may slide across the dividers. For instance, an element of $B_s \ot B_t$ consists of (a linear combination of) a choice of 3 polynomials (left, middle, and right), such that an
$s$-invariant polynomial may slide across the first divider, and a $t$-invariant polynomial across the second. An element of $R_w \ot R_v$ consists of 3 polynomials, and any polynomial may
be slid across any divider, at the cost of applying the appropriate element of $W$ to it. When we write $R_w$ in this way, right multiplication is the usual \emph{untwisted} multiplication
on the right slot; it is the left slot which is identified as a $\Bbbk$-module with $R$ in the definition of the previous paragraph. We call the element $1 \ot 1 \ot 1 \ot \cdots \ot 1$ of
such a bimodule a \emph{1-tensor}. Clearly the 1-tensor is the unique element of minimal degree, up to a scalar.

Let us assume for the rest of this discussion that the realization is faithful. We have
\[
\Hom(R_w,R_v) = \begin{cases} R & \text{if $w = v$} \\ 0 & \text{otherwise.} \end{cases}
\]
In other words, $\StdBim$ is isomorphic to the graded 2-groupoid for $W$ over $R$ (see \cite{EWFenn} for terminology). Therefore, any map between standard bimodules is determined by the image of the 1-tensor.  In particular, a degree $0$ map in $\StdBim$ between $R_\xb$ and $R_\yb$ will send the 1-tensor to a scalar multiple of the 1-tensor.

We write $\StdBim_0$ for the ungraded monoidal category whose objects are standard bimodules and whose morphisms are degree $0$ maps. This is equivalent as a monoidal category to the
2-groupoid for $W$ over $\Bbbk$, and can be defined without any restrictions on $\Bbbk$. Presenting this 2-groupoid as a monoidal category by generators and relations is surprisingly
interesting (see \cite{EWFenn} and the next chapter).

We note that all tensor products of Soergel bimodules and standard bimodules are free as left (or right) $R$-modules, and thus the Hom spaces between them are $R$-torsion-free. We will soon show that, under certain assumptions on the realization, they are free as $R$-modules.

Now we discuss the maps between Soergel bimodules and standard bimodules. There is an injection of bimodules $R(-1) \to B_s$ arising from the Frobenius algebra structure, sending $1
\mapsto \Delta_s$ (defined in the previous section). The cokernel of this map is naturally isomorphic to $R_s(1)$, via the map $B_s \to R_s(1)$ sending $f \ot g \mapsto fs(g)$. Conversely,
there is a surjection $B_s \to R(1)$ sending $f \ot g \mapsto fg$. The kernel of this map is naturally isomorphic to $R_s(-1)$ via the map sending $1 \mapsto \d \ot 1 - 1 \ot \d$.

\begin{remark} Just as $\Delta_s$ has a canonical description, so too does this element $\d \ot 1 - 1 \ot \d$. The trace map $\pa_s \co R \to R^s$ induces an $R$-bilinear pairing $(f,g)
\mapsto \pa_s(fg)$, but it also induces a twisted-bilinear pairing $(f,g) \mapsto \pa_s(f s(g))$. The bases $\{1,\d\}$ and $\{-\d,1\}$ of $R$ over $R^s$ are dual for the twisted pairing,
and the element $\d \ot 1 - 1 \ot \d$ is independent of the choice of dual bases. The reader can ponder the notion of a twisted Frobenius extension. \end{remark}

All four of these maps have graded degree $1$. We encode them in two short exact sequences: \begin{equation} 0 \to R(-1) \to B_s \to R_s(1) \to 0 \label{ses1} \end{equation}
\begin{equation} 0 \to R_s(-1) \to B_s \to R(1) \to 0 \label{ses2} \end{equation} Thus $B_s$ is filtered by $R$ and $R_s$, though in no particular order, and the grading shifts which
appear depend on the chosen order. This implies that every Bott-Samelson bimodule has a filtration whose subquotients are standard bimodules (with shifts). Because each bimodule $R_w$ is
indecomposable, any direct summand of a Bott-Samelson bimodule (and hence any Soergel bimodule) also has such a filtration. However, the order in which standard modules appear in such a
filtration need not respect the Bruhat order.

Now assume that $\hg$ is a Soergel realization. It is a deeper fact that any Soergel bimodule has a filtration in which all successive subquotients are standard modules, occurring in an
order refining the Bruhat order (resp. the reversed Bruhat order). Such filtrations are called \emph{standard filtrations}, and the
graded multiplicities of the standard modules appearing do not depend on the choice of filtration. One defines the \emph{character} of a Soergel
bimodule $B$ as the element $\ch(B)$ of the Hecke algebra counting these graded multiplicities. (For the precise definition of the character see \cite[\S 5]{Soe3}).

Soergel went on to prove that the space of homomorphisms between any two Soergel bimodules is free as a left or right $R$-module, with graded rank given by evaluating the standard pairing
on the characters of each bimodule. Soergel used this formula to classify, in a non-constructive way, all the indecomposable Soergel bimodules. From the above discussion it is clear that
if $\un{w}$ is a rex, then $B_{\un{w}}$ has $R_w$ in its standard filtration with multiplicity $1$, and all other standard modules appearing are isomorphic to $R_y$
for $y<w$. The following theorem is due to Soergel \cite[Satz 6.14]{Soe3}:

\begin{thm} Let $\hg$ be a Soergel realization of $W$. For all $w \in W$ there exists a unique (up to isomorphism) bimodule $B_w$ which occurs as a direct summand of $B_{\un{w}}$ for any reduced expression $\un{w}$ for $w$. The bimodule $B_w$ is uniquely determined as a summand of $B_{\un{w}}$ by the fact that it is indecomposable and (some shift of) $R_w$ occurs in its standard filtration. The set $\{ B_w \; | \; w \in W \}$ constitutes a complete set of non-isomorphic indecomposable Soergel bimodules, up to isomorphism and grading shift.\label{SoergelThm1} \end{thm}

The statement that $B_{\un{w}}$ contains a unique summand such that (some shift of) $R_w$ occurs in its standard filtration follows immediately from the indecomposability of $R_w$. The fact that such summands for different rexes are isomorphic is not difficult to prove (using the Krull-Schmidt property and an idempotent lifting argument). The difficulty in the theorem is to show that any summand of any Bott-Samelson bimodule is isomorphic to one of the bimodules $B_x$ up to a shift.

The theorem implies that $B_w$ is the only summand of $B_{\un{w}}$ which is not a summand of (some shift of) $B_{\un{y}}$ for any shorter sequence $\un{y}$. In principle one can
``construct" $B_w$ by finding all the summands of $B_{\un{w}}$ which occur as shifts of summands of lower terms, removing them, and seeing what remains. This amounts to a calculation of
all idempotents in $\End_0(B_{\un{w}})$, which is a difficult and subtle question.

Because of the implicit definition of the indecomposable bimodule $B_w$, its intrinsic properties often depend on the characteristic of $\Bbbk$ and the realization $\hg$ of $(W,S)$.
For example, it may happen that $B_{\un{w}}$ admits a non-trivial decomposition in characteristic $0$, with a nontrivial summand $B_w$, while for a certain finite characteristics $B_{\un{w}}$ is indecomposable, meaning that $B_w = B_{\un{w}}$. In this paper we determine the algebra of endomorphisms of $B_{\un{w}}$, but are only able to make very basic statements about the representation theory of this algebra.

%===========
\subsection{Categorification}
\label{subsec-SoergelCatfn}
%===========

We denote the split Grothendieck group of an additive category by $[\CC]$. That is, $[\CC]$ is the abelian group generated by symbols $[M]$ for all objects $M \in \CC$ subject
to the relations $[M] = [M'] + [M'']$ whenever $M \cong M' \oplus M''$ in $\CC$. When $\CC$ is monoidal, $[\CC]$ has the structure of a ring via $[M][M'] = [M\otimes M']$. If in addition
$\CC$ is graded with grading shift functor $M \mapsto M(1)$, then $[\CC]$ has the structure of a $\Zvv$-algebra via $v[M] \define [M(1)]$.

The following is \emph{Soergel's categorification theorem} \cite[Satz 1.10 and Satz 5.15]{Soe3}:

\begin{thm} Let $\hg$ be a Soergel realization of $W$. There is a unique isomorphism of $\Zvv$-algebras:
\begin{align*}
\e: \HB &\simto [\SBim] \\
\un{H}_s &\mapsto [B_s].
\end{align*}
Given Soergel bimodules $B$ and $B'$ the graded rank
of $\Hom(B,B')$ as a free left (or right) $R$-module is given by $(\e^{-1}[B],\e^{-1}[B'])$, where $(-,-)$ denotes the standard pairing on $\HB$.
\label{SoergelThm2}
\end{thm}

The uniqueness of $\e$ is immediate, because $\{\un{H}_s\}_{s \in S}$ generates $\HB$. To see that $\e$ is a homomorphism it is enough to check the relations \eqref{rank1hecke} and
\eqref{rank2hecke}. Using the Frobenius algebra structure and the isomorphism $R \cong R^s \oplus R^s(-2)$ of $R^s$-bimodules, one can easily check that \begin{equation} B_s \ot B_s \cong
B_s(1) \oplus B_s(-1). \label{rank1BS} \end{equation} This isomorphism categorifies equation \eqref{rank1hecke}. Under certain assumptions on the realization, the categorification of
equation \eqref{rank2hecke} comes from an explicit description of $B_w$ as the image of a certain idempotent, for every $w$ contained in a standard rank 2 parabolic subgroup. More details
can be found in \cite[\S 4]{Soe3} or \cite{EDihedral}. In this paper we will give an alternate proof of the categorification of \eqref{rank2hecke} in more generality. It follows that $\e$
is a homomorphism of $\Zvv$-algebras.

Once one knows that $\e$ is a homomorphism, the statement that it is an isomorphism follows from the  classification of indecomposable Soergel bimodules and the fact that their characters are upper triangular. More precisely, if we fix a rex $\un{w}$ for every $w \in W$ then the set $\{ H_{\un{w}} \; | \; w \in W\}$ is easily seen to be upper triangular in the standard basis of the Hecke algebra with respect to the Bruhat order, and hence is a basis. On the other hand, Theorem \ref{SoergelThm1} shows that $\{ [B_{\un{w}}] \; | \; w \in W \}$ is upper triangular in the basis $\{ [B_w] \; | \; w \in W \}$ for $[\SBim]$. As $\un{H}_{\un{w}}$ is mapped to $[B_{\un{w}}]$ it follows that $\e$ is an isomorphism.

In fact, Soergel shows that the character map
\[
\ch : [\SBim] \to \HB
\]
 discussed in previous section provides an inverse to $\e$. The character map is rather subtle. In general it is not known how to describe the element $\ch(B_w)$ in the Hecke algebra. However, if $\Bbbk$ is a field of characteristic zero Soergel proposed the following conjecture, which came to be known as \emph{Soergel's conjecture}:
 
\begin{conj} If $\Bbbk$ is a field of characteristic $0$ then $\ch(B_w) = \un{H}_w$. \end{conj}

It is immediate that Soergel's conjecture implies the Kazhdan-Lusztig positivity conjectures. Earlier work of Soergel \cite{Soe1} showed that his conjecture also implies the
Kazhdan-Lusztig conjecture. Soergel's conjecture has recently been proved by the authors in \cite{EWHodge} for realizations (not necessarily reflection faithful) defined over $\RM$, and satisfying a certain positivity property.

\begin{remark} This statement of Soergel's conjecture is known to fail when $\Bbbk$ is a field of finite characteristic $p$. The image of $[B_w]$ in $\HB$ is known as the $p$-\emph{canonical basis}, and its computation is an interesting open problem (see \cite{Wpcan}).
\end{remark}

\begin{remark} Soergel's conjecture is false for certain complex realizations of the affine Weyl group of type $A$ which do not admit any real form (as in part (6) of Example
\ref{ex:realizations}, when $q$ is a root of unity). See \cite{EQAGS} for details. Hence the assumption that the underlying representation be defined over $\RM$ is essential for Soergel's
conjecture to hold. (The authors' proof of Soergel's conjecture uses positivity considerations in a crucial way!) \end{remark}

The fact that $\un{H}_s$ is self-biadjoint in the standard pairing is categorified by the fact that the functor of tensoring with $B_s$ is its own left and right adjoint. Just as in Remark \ref{rmk:traces pairings}, any Hom space between Bott-Samelson bimodules can be understood after adjunction in terms of $\Hom(R,B_{\un{w}})$. Soergel's Categorification Theorem
combined with Corollary \ref{cor:defects} implies that this space is free as a left $R$-module with a homogeneous basis in bijection with the subsequences of $\un{w}$ with end-point $e$.
Following the work of Libedinsky \cite{Lib}, we will show that this basis can be constructed in a natural way.

\begin{remark} Suppose that $(W_1, S_1)$ and $(W_2, S_2)$ are Coxeter systems and $S_1 \subset S_2$ is an inclusion realizing $W_1$ as a standard parabolic subgroup of $W_2$. Suppose also
that we have a $W_1$-equivariant inclusion of realizations $\hg_1 \subset \hg_2$. There is an obvious functor $\iota \co \BSBim_{S_1} \to \BSBim_{S_2}$ sending $B^1_s \mapsto B^2_s$.
Soergel's Categorification Theorem implies that this functor is ``fully-faithful after base change.'' That is, if $R_1$ and $R_2$ denote the polynomial functions on $\hg_1$ and $\hg_2$
respectively then \[ R_2 \otimes_{R_1} \Hom(B^1_{\un{w}},B^1_{\un{w}'}) = \Hom(B^2_{\un{w}}, B^2_{\un{w}'}). \] Said another way, homomorphisms from $B_{\un{w}}$ to $B_{\un{w}'}$ only
depend on the simple reflections which occur in $\un{w}$ and $\un{w}'$ (up to the action of polynomials). The (diagrammatic) proof that these functors are fully faithful after base change
will be clear later in this paper. As a result, one may define the category of Soergel bimodules for $(W, S)$ even when $S$ is infinite, either by working with an infinite dimensional
realization of $W$, or by taking a limit over the categories associated to finite subsets of $S$.\label{includingindexsets} \end{remark}

%%%%%%%%%%%%%%%%%%%%%%%%%%%%%%%%%%%
\subsection{Localization}
%%%%%%%%%%%%%%%%%%%%%%%%%%%%%%%%%%%

As $\Bbbk$ is a domain, so is $R$. Let $Q$ denote the fraction field of $R$, which is an (ungraded) field living over $\Bbbk$. One can extend any $R$-module to a $Q$-module in the usual
way. We have already shown in general that any tensor product of standard and Soergel bimodules is free as a left (or right) $R$-module, and that morphism spaces are torsion-free.
Therefore, extension of scalars from $R$ to $Q$ on the left (or right) will yield a free $Q$-module, and will act faithfully on morphisms.

\begin{lem}\label{localizationononeside} For any $\un{w} = s_1s_2 \dots s_d$ one has an isomorphism of right $Q$-modules
\[
B_{\un{w}} \otimes_R Q = Q \otimes_{Q^{s_1}} Q \otimes_{Q^{s_2}}  \dots \otimes_{Q^{s_d}} Q.
\]
The analogous statement can be made for standard bimodules, or tensor products of standard and Bott-Samelson bimodules.
\end{lem}

\begin{proof} We show that $B_s \otimes_R Q \cong R \otimes_{R^s} Q$ is isomorphic to $Q \ot_{Q^s} Q$, which implies the lemma by a simple induction. One has an obvious inclusion $i : R \otimes_{R^s} Q \into Q \ot_{Q^s} Q$ and it is enough to show that this is surjective. However, any $0 \ne f \in R$ divides the $s$-invariant polynomial $fs(f)$ and in $Q \ot_{Q^s} Q$ we have
\[
s(f) \otimes \frac{1}{fs(f)} = \frac{1}{f} \otimes 1.
\]
It follows that $i$ is surjective, and the lemma follows. The analogous proof for standard bimodules is even simpler. \end{proof}

In particular, base change from $R$ to $Q$ on the right will send a Bott-Samelson bimodule to a $Q$-bimodule, not just an $(R,Q)$-bimodule, and the base change functor is monoidal. Let
$\BSBim_Q$ denote the essential image of $\BSBim$ inside the category of $Q$-bimodules. By the above lemma, it is the full subcategory of $Q$-bimodules generated by the bimodules
$B_{s,Q}\define Q \ot_{Q^s} Q$. Let $\SBim_Q$ denote its Karoubi envelope. We define $\StdBim_Q$ similarly.

As $Q$-bimodules we have a splitting \begin{equation} \label{eq:split} B_{s,Q} \cong Q \oplus Q_s. \end{equation} One way to observe this is that the sequences \eqref{ses1} and
\eqref{ses2} are split. In fact, they ``split each other:" the composition $R \to B_s \to R$ is multiplication by $\a_s$, which becomes invertible in $Q$. Therefore, any Soergel bimodule
over $Q$ splits into standard bimodules. When the realization is faithful one has \[ \Hom(Q_w,Q_v) = \begin{cases} Q & \text{if $w = v$} \\ 0 & \text{otherwise,} \end{cases} \] so that
standard bimodules over $Q$ are indecomposable and pairwise non-isomorphic. It follows that $\SBim_Q$ is equivalent to the (additive closure of the) 2-groupoid of $W$ over $Q$.

We conclude $\SBim_Q$ has a drastically different behaviour to $\SBim$ --- $\SBim$ is far from being a 2-groupoid! In the following, we will refer to the faithful monoidal functor
\begin{gather*}
\SBim \to \SBim_Q
\end{gather*}
as \emph{localization}. It plays an important role in what follows, essentially because $\SBim_Q$ is such a simple category. We already know that Hom spaces between various $Q_w$ are either rank 1 or 0. Because the passage to the localization is injective, and because morphisms between various $Q_w$ are so easy,
we can study morphisms in $\SBim$ based on what their restrictions are to simple summands in $\SBim_Q$ after localization. Given an explicit morphism $f$ in $\BSBim$, we will be able to
write down explicitly its restriction to each standard summand of the source and target in $\SBim_Q$, and in doing so, will be able to conclude whether or not $f=0$.

\begin{remark} One may use the homogeneous fraction field $Q'$ instead of the full fraction field, and the results in this paper pertaining to localization will be (essentially) unchanged.
The homogeneous fraction field is graded; however, $Q' \cong Q'(2)$ so that the grading does not yield any interesting invariants (even the parity is irrelevant, as investigation reveals).
Though more difficult, all the localization results in this paper should hold if one only inverts the simple roots $\a_s$ and their $W$-conjugates. One advantage to using
$Q$ over these alternatives is that $Q$ is a field. \end{remark}

% %===========
% \subsection{Extension of scalars}
% \label{subsec-localandleft}
% %===========
% 
% {\color{green} FIX THIS SECTION}
% 
% Localization is an special case of extension of scalars. Although localization is the only example of extension of scalars used in this paper, we feel that the reader should know
% about the more general setting.
% 
% Forgetting the monoidal structure on Soergel bimodules allows us to view $\SBim$ as an additive category enriched in graded right $R$-modules. (We arbitrarily decide to view the homomorphisms spaces between Soergel bimodules as (free) right $R$-modules. It would change little below to consider the homomorphism spaces as left $R$-modules instead. The fact that these homomorphism spaces are free is important, which explains why we do not choose to view them as $R \ot_{\Bbbk} R$-modules.) Given any (graded)  $R$-algebra $R'$ we obtain a new (graded) additive category $\SBim_{R'}$ by tensoring all Hom spaces on the right by $R'$ and taking the Karoubi envelope. In this fashion one can hope to understand Soergel bimodules as a ``family of categories over $\Spec R$''. The fact that all homomorphism spaces are free tells us that this family is flat. This philosophy plays a key role when relating categories of Soergel bimodules to representation theory (usually going by the name ``deformation'').
% 
% For any such $R'$ one has a natural morphism
% \begin{equation} \label{eq:extendcompare}
% \Hom_{(R,R)}(B, B') \otimes_R R' \to \Hom_{(R,Q)}(B\otimes_R R', B' \otimes_R R').
% \end{equation}
% which is an isomorphism if $R'$ is a (homogenous or inhomogenous) localization of $R$.\footnote{GW: Still need to check that this is OK.}In particular, if $R' = Q$ we recover the localization functor of the previous section. (This also explains why the notation $\SBim_Q$ is not ambiguous.) Note that in this description it is immediately obvious that this functor is monoidal.
% 
% Another important example is obtained by considering the $R$-algebra $\RM = R/ R^+$, where $R^+ \subset R$ denotes the ideal of elements of positive degree. In this case $\SBim_{\RM}$ is no-longer a monoidal category. By the graded Nakayama lemma much of the complicated structure of $\SBim$ is still present in $\SBim_{\RM}$. In this setting, we do not know whether \eqref{eq:extendcompare} is an isomorphism in general. For finite Weyl groups this holds by \cite[Proposition 8]{Soe2}.
% 
% \begin{remark}
% In \cite{Soe1}, Soergel shows that if $W$ is a finite Weyl group then $\SBim_{\RM}$ is a graded version of category $\OC$:
% if we set $B\define \bigoplus_{x \in W} B_x$ then then the category of (ungraded) modules over the endomorphism algebra
% \[
% A = \End_{\SBim_{\RM}} (B) = \End_{R}( \oplus_{x \in W} B_x \otimes_R \RM)
% \]
% is equivalent to the principal block of category $\OC$, for any complex semi-simple Lie algebra with Weyl group $W$. Surely the study of the corresponding categories for non-crystallographic $W$ holds some surprises!
% \end{remark}
% 
% 
% 
% 
% 
% {\color{red} COMMENTS:
%   \begin{enumerate}
%   \item In opposition to the left quotient, which contains all the interesting homological features of category $\OC$, the localization to $Q$ is semi-simple. {\color{blue} This is false because $Q$ is not a semi-simple ring. It would be true if we passed to the full fraction field of $R$. {\color{red} (Which we have done now)}}
%   \item 
% First we consider the \emph{left quotient} of $\SBim$, which is the subcategory of right $R$-modules obtained by taking all Soergel bimodules and setting the left action of $R_+$ to
% zero. In this category, all Hom spaces are finite-dimensional (over $\Bbbk$), with dimension given by (a rescaling of) the standard pairing. {\color{blue} As far as I am aware it is not known whether this holds outside of Weyl groups, and the proof for Weyl groups is subtle (using deformations of projectives). I would like to know whether it holds in general, but it is not on top of my list of things I would like to know:)}
%   \end{enumerate}
% 
% WOULD DELETE: First we consider the \emph{left quotient} of $\SBim$, which is the subcategory of right $R$-modules obtained by taking all Soergel bimodules and setting the left action of $R_+$ to
% zero. In this category, all Hom spaces are finite-dimensional (over $\Bbbk$), with dimension given by (a rescaling of) the standard pairing. This category is equivalent to the
% additive category of projectives in category $\OC$, in case $W$ is a Weyl group. Because any symmetric polynomial acts the same on the left and right of any Soergel bimodule, we may
% as well assume that we are dealing with right $C$-modules. The left quotient contains roughly all the complexity of the original category $\SBim$. It will not be of much relevance
% in this paper, but we include it so the reader can taste more of the connections.}
% 
